% ===== 附录 =====
\section{关键参数与输出说明}
\label{sec:r-appendix}

为保证评委版正文保持论文式叙事，本附录集中列出当前演示版本中最关键的复现参数和输出结构。这里记录的是本文结果所对应的固定设置，不应与未来可能重新校准的参数混淆。

\begin{table}[h]
\centering
\caption{当前评委版结果对应的关键参数。}
\label{tab:r-appendix-params}
\footnotesize
\renewcommand{\arraystretch}{1.12}
\begin{tabular}{>{\raggedright\arraybackslash}p{4.0cm}>{\raggedright\arraybackslash}p{4.0cm}>{\raggedright\arraybackslash}p{4.6cm}}
\toprule
项目 & 当前设置 & 说明 \\
\midrule
实例分割模型 & \texttt{IG2\_full\_set\_cp\_SAM} & ImageGrains 2.0 / Cellpose-SAM 预训练模型 \\
演示图尺度 & 0.208~mm/px & 来自现有样本产物中的已用尺度 \\
筛分等效参数 & $\boldsymbol{\theta}=(1,0,0)$ & 当前以短轴 $b$ 为未校准基线 \\
质量代理指数 & $\gamma=3$ & 几何相似条件下的体积/质量尺度先验 \\
筛级边界 & 5, 10, 16, 20, 25, 31.5, 40~mm & 同时保留 $<5$ 与 $\ge40$~mm 边界区间 \\
针片状候选 & $AR<0.50$ & 二维投影规则，不等价于三维量规真值 \\
圆形候选 & $C>0.85,S>0.95,AR>0.80$ & 几何阈值规则 \\
棱角状候选 & $(C<0.60\ \text{or}\ S<0.90)$ 且非针片状 & 几何阈值规则 \\
超限异常 & $d>50$~mm & 尺寸型异常规则 \\
疑似泥团 & $40<d\le50$~mm 且 $S<0.85$ & 可关闭的保守几何候选规则 \\
\bottomrule
\end{tabular}
\end{table}

每个场景的最小复现结果包包括：原始图像与尺度信息、integer instance mask、实例/轴长可视化、逐颗粒 CSV、场景级 summary、级配曲线以及形貌/异常可视化。逐颗粒表保存长短轴、面积、周长、等效直径、筛分等效粒径、质量代理权重和规则标签，因此场景级 $D_p$ 与筛级占比可以在不重新执行分割模型的情况下重新计算。

本文的三张演示样本为 \texttt{agg\_001}、\texttt{agg\_005} 和 \texttt{agg\_029}。评委版正文报告的所有实例数、$D_{10}/D_{50}/D_{90}$、形貌比例和异常数量均来自这三组现有产物；没有使用模型生成的“伪真值”补充外部准确率。